\documentclass[platex,dvipdfmx,12pt]{jsarticle}
%使ってるのがplatexだからクラス指定はplatex.間違えてもuplatexにしないように

\usepackage{graphicx}
\usepackage[hyphens]{url}
\usepackage{otf}
\usepackage{multirow}w
\usepackage{amsmath}
\usepackage{here}
\usepackage{geometry}
\usepackage{ascmac}

\usepackage[hang,small,bf]{caption}
\usepackage[subrefformat=parens]{subcaption}
\captionsetup{compatibility=false}
\usepackage{textcomp}

\usepackage{listings} %日本語のコメントアウトをする場合jvlisting（もしくはjlisting）が必要
%ここからソースコードの表示に関する設定
\lstset{
  basicstyle={\ttfamily},              % 基本:等幅フォント
  identifierstyle={\small},            % 識別子
  commentstyle={\smallitshape},        % コメント:斜体
  keywordstyle={\small\bfseries},     % キーワード:太字
  stringstyle={\small\ttfamily},      % 文字列
  frame={tb},                           % 上部と下部に線を表示
  breaklines=true,                      % 行が長い場合に折り返す
  columns=[l]{fullflexible},            % 列幅を自動調整する（見た目が良くなる）
  numbers=left,                         % 行番号を左側に表示
  xrightmargin=0zw,                     % 右マージンのサイズ．
  xleftmargin=1.6zw,                    % 左マージンのサイズ．行番号が2桁でも行左端からはみ出ない値．
  numberstyle={\scriptsize},           % 行番号のスタイル．スクリプトサイズのフォントを使用
  stepnumber=1,                         % 何行ごとに行番号を表示するか
  numbersep=1zw,                        % 行番号とソースコードの間の距離
  lineskip=-1.4ex,                      % ソースコードの行間（結構詰め気味）  
  morestring=[b]",                     %日本語等のStringを正常にダブルクオーテーションで囲む
  showstringspaces=false,               %Stringに囲まれた部分ではスペースがSymbol付きで表示されるのでそれを消す
}
\renewcommand{\lstlistingname}{ソースコード}

\geometry{left=20mm,right=20mm,top=20mm,bottom=30mm}
\renewcommand{\baselinestretch}{1.2}

\makeatletter %sectionをchapterっぽくする.
\renewcommand{\presectionname}{第}
\renewcommand{\postsectionname}{章}%はコメントですがつけ忘れると\section*でバグります

% 図番号を1.5のようにする
\renewcommand{\theequation}{\arabic{section}.\arabic{equation}}
\renewcommand{\thefigure}{\arabic{section}.\arabic{figure}}
\renewcommand{\thetable}{\arabic{section}.\arabic{table}}

\renewcommand{\section}{%
\clearpage
\@startsection{section}{1}{\z@}%
{\Cvs \@plus.5\Cdp \@minus.2\Cdp}% 見出し前に付けられる空白量
{.5\Cvs \@plus.3\Cdp}% 見出し後に付けられる垂直方向の空白量
{\normalfont\LARGE\headfont\raggedright}}% 見出しのスタイル
\makeatother


\begin{document}

\begin{titlepage}
\begin{center}
\begin{huge}
情報工学実験 \ajRoman{2} 第6回レポート \linebreak \vspace{1em}
GUI開発とソフトウェアテストに関する実験\end{huge}
\end{center}

\vspace{6em}

\begin{LARGE}
\begin{flushright}
4年情報工学科 34番 \linebreak \vspace{1em}
増井 崇
\end{flushright}
\end{LARGE}

\vspace{4em}

\begin{Large}
\begin{tabular}{ll}
提出期限 : & 2026年01月30日（金）13:00 \\
提出日 : & 2026年01月27日（火）23:59
\end{tabular}

\vspace{2em}

共同実験者：5班

\vspace{1em}

\begin{tabular}{ll}
29番 & 中村 奏月 \\
30番 & 西陰寺 厚志 \\
31番 & 坂東 惠佑 \\
32番 & 福田 悠人 \\
33番 & 前田 貴亮 \\
35番 & 増本 稜大 \\
\end{tabular}

\end{Large}

\end{titlepage}

\section*{アブストラクト}
\thispagestyle{empty}
本レポートでは，設計書を用いて，WinFormsを使ったGUIプログラムの開発を行い，設計書に準拠したプログラムのプログラミング開発や，プログラムの動作テストなどを設計書を元にMSTestsなどを用いて作成し，実行した．
設計書を用いたGUIプログラムの開発では，小規模な四則演算プログラム，中規模のマルバツゲームプログラム(単純AI，Q学習AI)，大規模のオセロゲームプログラムを
開発した．
特に，四則演算およびマルバツゲームプログラムはスクラッチ開発であり，仕様書に書かれた要件より，ゼロベースでクラス設計やメソッド設計など行い，実装まで行なった．
さらに，マルバツゲームプログラムのプログラミングにおいては，以前に市川先生の実験で作成したマルバツゲーム専用AIとの対戦機能の実装などを行なった．
その際には，上野先生のマルバツゲームの要件と，市川先生の実験で作成したマルバツゲーム環境の要件との差異を吸収するためのラッパーメソッドの実装を行なった．
また，オセロゲームの開発では，事前に用意された多くのプログラム群に対して，追加で必須となるメソッドおよびクラスを実装するという，
既存プログラムの改修として実装を行なった．
特に，CalcConditionやClickedメソッドなど，設計書の動作要件が複雑なプログラムを，既存のメソッドを使いつつ，
Privateメンバへのアクセス制限などを回避し，規定の仕様を満たすプログラムの開発を行なった．
最終的に，作成したすべてのプログラムにおいてテストを行なった．
四則演算プログラムおよび，マルバツゲームプログラムでは，GUI上で，指定されたテストケースのテストを行い，そのテストケースの意味を考察し，テストケースの十分性をレポート内で示した．
また，オセロゲームのプログラムでは，BoardクラスのCalcConditionメソッドやCellのコンストラクタ，Setメソッドを
MSTestのTestMethodを用いて，機械的に大量のテストを行なった．
テストケース設計においては，テスト対象の入力，動作要件から同値分割法や境界値分析，デシジョンテーブルなどを用いて，
盤面上の膨大な状態数を，いくつかのクラスに分け，その中からテストすることでテストケースを削減することに成功した．
まとめでは，CI/CDとTDDがどのように組み合わさり，プログラムに対するテストの重要性を考察した．


\section{実験目的}
\setcounter{page}{1}
本実験では，仕様書を用いた実践的なソフトウェアプログラミングを行い，
実際の成果物に対してテストを行うなどの一連の流れを体験し，問題点や注意点，テスト手法を得ることを目的としている．

特に，仕様書を用いたプログラミングでは，GUIアプリケーションを題材に，C\#を用いて，イベント駆動プログラミングを実践し，実際のプログラミングに近い環境を再現し，経験することを目的としている．

また，テストでは，仕様書から動作要件を洗い出し，あるメソッドの作用を境界値分析や同値分析を用いて効率的に，自らが実装したプログラムの動作確認を体験することを目的としている．

最終的に，CI/CDやテスト駆動プログラミングなどのモダナイゼーションされたプログラミング環境の入口となるような実験をすることを目指す．


\section{実験原理}
\subsection{MVCモデル}
MVCモデルとは，アプリケーションの基本機能をModel，View，Controllerの3種類に分類し，各機能の責任を限定的にするための，GUIプログラミングにベストプラクティスの一つである．

特に，WindowsアプリケーションやJavaアプリケーションとの相性がよく，現在に至るまで，多くのソフトウェアプログラミングにおいて採用されてきた手法である．

\paragraph{Modelの働き}
Modelはデータモデルやリポジトリーパターンなどが実装されたクラスが属する分類である．
実際には，ViewやControllerが利用したいデータの，流れや型などを定義し，安全に上位層でデータを利用するための基盤を提供する役割を担う．

具体的には，DBやAPI，WebSocketなど多岐にわたる通信方法による差異をなくし，一定のデータクラスに格納することで，型安全性などを確保しつつ，データの流れを明確にできるなどの効果がある．
また，データの流れが明確になることで，モックデータとの差し替えが行いやすくなり，実際のプログラミングにおいて，テスト作業の簡略化やプログラム上でのテストのしやすさの向上などに寄与することもある．

\paragraph{View}
Viewは実際にエンドユーザーに閲覧される部分を構成するクラスが属する部分である．
このクラスでは，GUIアプリケーションにおける，ボタン要素であったり，リスト要素，HTML要素などをレンダリングし，見た目を宣言するなどの役割を担う．

具体的には，ボタン配置であったり，リスト要素の中身などをプログラムとして宣言し，それらを配置し，描画エンジン(SwingやWinForms)に渡すなどの作業を行う．

\paragraph{Controller}
Controllerでは，エンドユーザーのView上でのアクションを直接ハンドリングし，具体的な処理を施すクラスなどが属する．
実際には，View上でのクリックイベントや更新にかかるローディングイベントなどをハンドリングし，正常に入力された情報や触覚情報などを処理する役割を担う．

具体的には，ゲームの進行やModelクラスのリポジトリなどからデータをフェッチし，Viewへ渡すなどの作業を行う．


ただし，近年のWeb中心の設計（React，Flutter）およびAndroidなどのスマートフォンアプリプログラミング（Kotlin）では，ControllerとViewを分離することが難しくなっている．
具体的には，ViewがControllerに依存しすぎるためである（状態管理ライブラリなどによる）．
そのため，Viewは状態を持たず，ControllerがViewの状態も保持するように設計を進める，MVVMモデルというものも登場している．

MVVMモデルでは，Model, ViewModel, Viewの三要素が状態という概念で密接に結びつく，プログラミングモデルである．

MVCとの違いとしては，ViewModelがViewの状態を保持し，Viewはただ，描画することだけに専念するという完全分離型の設計なことが挙げられる．
ViewとViewModelはObserverパターンというデザインパターンを使って，即時データバインディングが行われ，GUIにデータの更新が反映される．

両者の特徴を踏まえ，現在も静的な動き中心のGUIアプリケーションはMVCモデルに基づき，動的なリアルタイム性の求められるアプリケーションではMVVMモデルが採用されるようになっている
（本文における動的なリアルタイム性とは，地図アプリであったり，SNSなどのリアルタイムにサーバーやクライアントの情報更新をアプリケーションへ反映することが求められるような性質を指す）．


\subsection{デザインパターン}
デザインパターンとは，オブジェクト指向プログラミングの考え方に基づき，各クラスの責任の所在を明らかにしつつ，プログラム自体の利便性を向上させるための設計手法である．

デザインパターンにおけるデファクトスタンダードとして，GoFというデザインパターンをまとめた書籍から引用された23種のデザインパターンが存在する．
まとめると，次の表\ref{tab:gof-design-pattern}のようになる．

\begin{table}[H]
  \centering
  \caption{GoFと呼ばれる23種のデザインパターン}
  \label{tab:gof-design-pattern}
  \begin{tabular}{cc}
    % ここにカラムを入力
    名称 & 名称 \\
    \hline
    % ここに各セルのデータを入力
    Abstract         & Chain of Responsibility \\
    Builder          & Command                 \\
    Factory Method   & Interpreter             \\
    Prototype        & Iterator                \\
    Singleton        & Mediator                \\
    Adapter          & Memento                 \\
    Bridge           & Observer                \\
    Composite        & State                   \\
    Decorator        & Strategy                \\
    Facade           & Template Method         \\
    Flyweight        & Visitor                 \\
  \end{tabular}
\end{table}

特に，AbstractパターンやObserverパターンFactory MethodパターンやSingletonパターンなどは非常によく用いられる．
AbstractパターンはAbstractクラスとして，
ObserverパターンはNotifierなどとして，
Singletonパターンは認証情報や画面描画のインスタンス取得の際の概念として，
Factory Methodパターンはクラスのコンストラクタとして利用されることが多い．


\subsection{ソフトウェアテスト}
ソフトウェアテストは一般に以下の目的で行われる．

\begin{enumerate}
  \item バグの発見
  \item ソフトウェアの動作と仕様書の一貫性確認
  \item ユーザー要件の達成度の確認
  \item ソフトウェアの品質保証
\end{enumerate}

実際には，様々なテストが多段的に行われ，実使用のシナリオや，サイバー攻撃に寄せたシナリオ，
実使用から外れたシナリオなど，幾つものパターンを用いてテストを行う．

様々なテストが行われるものの，一般に以下の表\ref{tab:tests}に挙げる，4つのテストを段階的に行う．

\begin{table}[H]
  \centering
  \caption{一般的なテストの種類}
  \label{tab:tests}
  \begin{tabular}{ccc}
    % ここにカラムを入力
    テストの分類 & テスト対象 & テスト内容 \\
    \hline
    % ここに各セルのデータを入力
    ユニットテスト & コンポーネント & コンポーネントの入出力や作用を確認 \\
    結合テスト & 複数のコンポーネント & 正しいユニット同士の相互作用を確認 \\
    システムテスト & ソフトウェア全体 & 要件定義に従っているかの確認 \\
    受け入れテスト & ステークホルダ全体 & 発注元で正しく動作するか確認 \\
  \end{tabular}
\end{table}

コンポーネントとは，システムの機能を持つ最小単位であるため，具体的にはメソッドなどのことを指す．

実際に，プログラム主体のテストはユニットテストや結合テスト，システムテストまでで，受け入れテストでは，具体的な運用フローを含めて，クライアントの方々が使用できるかといった，
実地的な点までテストを行う．

単体テストからテストを開始し，動作を保証できるコンポーネントを増やしていくことで，最終的に大規模なコンポーネントの塊であるシステムのテストを完了するというルートで
テストを行うのが基本である．

\subsection{ユニットテストで用いられる手法}
ユニットテストをもう少し深く考えたい．

そもそも，ユニットテストでテストしたい対象とは，コンポーネントの動作であり，コンポーネントが最終的に出力する値や作用するオブジェクトの値を見ることで動作を確認できる．
また，その動作は一般に，入力の組やコンポーネントの状態によって決定する．

したがって，入力の内容からコンポーネントの動作を，仕様書を用いて推測し，その内容と実際の動作の差を見ることで，コンポーネントの動作が意図したものであるのかを判断することができる．

この時，全ての入力と動作の組を確認することは一般に容易ではない．

そのため，ブラックボックステストやホワイトボックステストなど，ある一定の点に着目することで，テストに必要なケースの数を削減する手法が用いられる．

\subsubsection{ホワイトボックステスト}
本テストでは，コンポーネントの中身に着目して動作を検証する．
具体的には，実行する命令の種類や，分岐の発生や，条件分岐のデシジョンテーブルなどに着目し，それぞれを網羅する形で動作を検証する．

このテスト手法では以下の主に4つの手法を用いて，テストケースを作成する．

\begin{enumerate}
  \item 命令網羅
  \item 分岐網羅
  \item 条件網羅
  \item 複数条件網羅
\end{enumerate}

\paragraph{命令網羅}
実行する様々な種類の命令を，全て通るようなプログラムパスを考え，網羅していく手法．
具体的な，条件分岐などは考慮せず，ただ命令を実行することだけに注力する．

\paragraph{分岐網羅}
判定条件などにより，Yes, Noに分かれるようなプログラムであれば，
そのどちらものルートを全て通れるようにプログラムパスを考え，網羅する手法．

\paragraph{条件網羅}
条件の内容を考慮し，デシジョンテーブルに従い，条件が出力する内容自体に着目し，網羅する手法．

\paragraph{複数条件網羅}
条件網羅が，複数加わった場合に同様の手順でテストケースを網羅する手法．



\subsubsection{ブラックボックステスト}
本テストでは，コンポーネントの入力と動作にのみ注目し，検証する．
具体的には，コンポーネントの設計書を用いて，それらをクラスタリングし，その中から代表値を抜き出して動作を検証することで，代表値が属しているクラス全ての値に対して
テストをしたことと同義だと判断し，テストケースを簡略化する．

そのための手法として，主に以下の2つを利用する．

\begin{enumerate}
  \item 同値分割法
  \item 境界値分析
\end{enumerate}

\paragraph{同値分割法}
同値分割法では，以下の手法でテストを行う．
仕様書より，入力と動作の組を列挙し，同じ動作が期待される入力を一つのクラスとして分け，他全ての組みに対して同様の手法を適用する．
これにより，入力自体をいくつかのクラスに分ける．
次に，各クラスが正常な動作を行うべきなのかという，入力の要件から，あるクラスが有効なクラスなのか，無効なクラスなのかを分類する．
最後に，各クラスから代表値を出し，それに対する動作を検証することで，プログラムの動作が正常であるかを判断する．
特に，同じ動作になる入力のクラスを同値クラスと言い，有効な同値クラスを有効同値クラス，無効な同値クラスを無効同値クラスという．

\paragraph{境界値分析}
この分析手法では，一般にコンポーネントは，入力要件の上限値や下限値，同値分割法のクラス間の境界となっている入力値付近で，
バグを引き起こしやすいという性質を使ってテストを行う．

考え方としては，ホワイトボックステストの条件網羅が近い．

具体的には，各同値クラスを一定の順序で並べた際に，隣のクラスとの境界値となる値を抽出し，それらの値に対する動作をテストするという手法である．

境界値というバグを引き起こしやすい付近の動作をチェックすると同時に，同値クラスの値の正常性も同時に証明できるハイブリッドな手法である．

しかし，問題点として，入力が複雑であるほど，爆発的にテストケースが増加することにある．

そのため，同値分割法は入力値があまり複雑ではない，または，同値クラスが入力値全体に及んでいる場合などは同値分割法を用いるべきである．
しかし，その判定であれば，境界値分析を用いて詳細に分析する方が安全である．


\section{実験手法}
\subsection{演習1 四則演算アプリのテスト}
本演習では，WinFormsを用いて，四則演算アプリを作成し，その動作を検証する．

指定されたアプリの要件は次の通りである．

\begin{enumerate}
  \item 入力は画面上の2つのTextBox(値1，値2)を用いる
  \item 出力は画面上の1つの編集不可設定のTextBox(出力値)を用いる
  \item 計算種別はコンボリスト(加減算，乗除算)を用いて選択
  \item 計算ボタンを押すと，入力された2つの値を，指定された計算種別で計算し，結果が出力される
  \item 計算結果はDouble型でとし，出力される内容はDouble.toStringの文字列をそのまま出力する
  \item また，数値以外の入力や計算種別が未選択の場合は"エラー!"を出力する
\end{enumerate}

従って，上記のアプリは次の計算イメージを実行することになる．

\begin{align*}
  値1 {計算種別} 値2 = 出力
\end{align*}

本アプリをプログラミング後，以下の表\ref{tab:enshu1-tds}に示すテストデータ(以下，TD)を実行し，その内容を記録する．

\begin{table}[H]
  \centering
  \caption{演習1 テストデータ}
  \label{tab:enshu1-tds}
  \begin{tabular}{ccccc}
    % ここにカラムを入力
    ID & 値1 & 値2 & 計算方法 & 期待される出力 \\
    \hline
    % ここに各セルのデータを入力
    TD-01 & 3 & 4 & 足し算 & 7 \\
    TD-02 & 3 & 4 & 引き算 & -1 \\
    TD-03 & 3 & 4 & かけ算 & 12\\
    TD-04 & 3 & 4 & 割り算 & 0.75 \\
    TD-05 & a & 4 & 割り算 & エラー! \\
    TD-06 & 3 & a & 割り算 & エラー! \\
    TD-07 & 3 & 4 & [空白] & エラー! \\
  \end{tabular}
\end{table}

TD-05, 06, 07はTD-04がより，割り算のロジックが正常に動作することを基準として考え，
TD-04から変更された部分を一箇所に絞りつつ，値を変化させることで，"エラー!"を出力する原因を明らかにしている．

他の四則演算のテストケースに関しても同様である．


\subsection{演習2 マルバツゲームのロジック実装，および拡張}
本演習では，WinFormsを用いて，マルバツゲームのロジックを実装し，指定された要件を満たすアプリケーションをプログラミングする．
また，自らプログラミングしたアプリケーションに対して，指定のテストケースを実行し，動作の正確性を検証する．

マルバツゲームの要件としては次のとおりである．

\begin{enumerate}
  \item 画面上に9つのButtonオブジェクトを配置
  \item プレイヤーとCPUの重複せず，必ず手番が回るターン制のゲームとする
  \item 勝敗判定ロジックを含む
  \item 勝敗判定の結果はLabelなどを用いてGUI上に表示する
  \item CPUの手番では，左上から右方向に探索し，押せる場所で押すというロジックを採用する
\end{enumerate}

本アプリをプログラミング後，以下の表\ref{tab:enshu2-tds}に示すTDを実行し，その内容を記録する．

\begin{table}[H]
  \centering
  \caption{演習2 テストデータ}
  \label{tab:enshu2-tds}
  \begin{tabular}{ccc}
    % ここにカラムを入力
    ID & タイミング & 期待されるゲーム結果 \\
    \hline
    % ここに各セルのデータを入力
    TD-01 & 9手目がクリックされる前まで & プレイヤーの勝利 \\ 
    TD-02 & 9手目がクリックされる前まで & CPUの勝利 \\
    TD-03 & 9手目がクリックされた瞬間 & プレイヤーの勝利 \\
    TD-04 & 9手目がクリックされた瞬間 & 引き分け \\
  \end{tabular}
\end{table}


\subsection{演習3 オセロゲームの実装，およびオセロゲームのテスト}
本演習では，WinFormsを用いて，設計済みのオセロゲームをプログラミングし，動作テストを行う．
設計済みのオセロゲームには，複数の実態のないメソッドが存在するため，それらを実装し，正常に動作するように修正し，完成を目指す．

動作テストでは，複雑なパターンに関して，ここではテストせず，正常に勝敗判定ロジックなどの新しくプログラミングしたコンポーネントが動作しているのかどうかについて，検証していく．

新規にプログラミングする部分は以下の三箇所である．

\begin{enumerate}
  \item Cellクラス
  \item GameControllerクラスのClickedメソッド
  \item BoardクラスのCalcConditionメソッド
\end{enumerate}

また，最終的に作成するオセロゲームの要件としては次のとおりである．

\begin{enumerate}
  \item ルールは一般的なオセロに準ずる
  \item プレイヤーの勝敗判定，引き分け判定を正常に処理できること
  \item プレイヤーやAIのターン飛ばしが発生した場合は，どちらかにクリックさせないようにロックする(描画タイミングや内容はサンプルに準拠)
\end{enumerate}

オセロゲームをプログラミング後，以下の表\ref{tab:enshu3-tds}に示すTDを実行し，その内容を記録する．

\begin{table}[H]
  \centering
  \caption{演習3 テストデータ}
  \label{tab:enshu3-tds}
  \begin{tabular}{ccc}
    % ここにカラムを入力
    ID & 盤面の状態 & 期待される状態 \\
    \hline
    % ここに各セルのデータを入力 
    TD-01 & 全て石で埋まっている状態 & 勝敗が表示され操作を受け付けない \\
    TD-02 & 盤面に空きがあるが，全てのピースが黒 & 黒の勝利 \\
    TD-03 & 盤面に空きがあるが，全てのピースが白 & 白の勝利 \\
  \end{tabular}
\end{table}

各テストケースでは，必ず一ターン以上の経過を含むこととする．
また，TD-02, 03に関して，初期状態からその盤面を再現することは難しいため，GameControllerをテスト用に改変し，任意の盤面を作成し，テストを行なっても良いものとする．

これは，必ずClickedメソッドの動作を実行させ，意図された動作をしているかを追加で検証するためである．



\subsection{演習4 ユニットテストのテストケースの制作方法}
本演習では，ユニットテストの中でもブラックボックステストに属する境界値分析を用いて，指定されたメソッドの動作テストを行う．

指定されたメソッドは設計書より次のとおりである\cite{journal:sekkeisyo}．

\begin{quote}
string grade(int score, int abs) \\

score（成績）に応じて成績判定（右表参照）を返す．
ただし，abs（欠課時数）が15を超えたらscoreにか
かわらず”×”を返す．scoreは0～100，absは0～30
のみを受け付け，それ以外の入力があった場合，”
エラー”を返す．
\end{quote}

また，テスト対象メソッドの出力は次の表\ref{tab:enshu4-target-method-1}, \ref{tab:enshu4-target-method-2}の通りである．
$x$を点数入力, $y$を欠席入力とする．

\begin{table}[H]
  \centering
  \caption{テスト対象の入出力表 点数}
  \label{tab:enshu4-target-method-1}
  \begin{tabular}{cc}
    % ここにカラムを入力
    入力 & 出力 \\
    \hline
    % ここに各セルのデータを入力
    $80 \leq x \leq 100$ & 優 \\
    $65 \leq x \leq 79$ & 良 \\
    $60 \leq x \leq 64$ & 可 \\
    $40 \leq x \leq 59$ & 不可A \\
    $0 \leq x \leq 39$ & 不可B \\
  \end{tabular}
\end{table}

\begin{table}[H]
  \centering
  \caption{テスト対象の入出力表 欠席}
  \label{tab:enshu4-target-method-2}
  \begin{tabular}{cc}
    % ここにカラムを入力
    入力 & 出力 \\
    \hline
    % ここに各セルのデータを入力
    $15 \leq y \leq 30$ & スコアに関わらず x(バツ)\\
    $0 \leq y \leq 14 $ & スコアの内容に依存 \\
  \end{tabular}
\end{table}

以上を踏まえて，本メソッドにおける同値クラスと，境界値を示し，計算過程や，実際に作成されたテストケースが明確になるように示す．


\subsection{演習5 自前の実装プログラムのテスト}
本演習では，オセロゲームの完成に向けて制作した，コンポーネント(Cellクラス コンストラクタ，Cell.Set，Board.GameController)のテストを行う．

まず，CellクラスのコンストラクタおよびSetメソッドは，明確な実装指標が明記されているため，そこから入力値と出力値を洗い出し，同値クラスへの分割を行い，境界値に着目するならば，
効率よくテストケースを列挙することができると予想される．

反対に，BoardクラスのGameControllerはオセロの盤面に対する評価のメソッドであることしか明記されていない．
そのため，テストケースの最適化を行わない場合，オセロの盤面の状態数の分だけテストを行い，その度に全ての状態を確認する必要がある．
具体的には，$3^{64}$ケースのテストである．
これは，現実的ではないため，盤面の各状態や，仕様書からGameControllerの作用を洗い出し，その作用に注目し，同値クラスや境界値を見つけ，効率的にテストを行うことが重要である．


また，これらのコンポーネントをテストする際に，コンポーネント内で使用されている既存のメソッドやクラスの動作は全て正しいものとしてテストを進めて良いものとする．


\section{実験結果}
\subsection{演習1 四則演算アプリのテスト}
実際に表\ref{tab:enshu1-tds}に従い，テストした結果が以下の図\ref{fig:enshu1-result-1}, \ref{fig:enshu1-result-2}, \ref{fig:enshu1-result-3}, である．

\begin{figure}[H]
  \centering
  \begin{minipage}[t]{0.30\textwidth}
    \centering
    \includegraphics[width=.8\textwidth]{assets/enshu_1/id1.png}
    \subcaption{TD-01}
  \end{minipage}
  \hspace{0.01\textwidth}
  \begin{minipage}[t]{0.30\textwidth}
    \centering
    \includegraphics[width=.8\textwidth]{assets/enshu_1/id2.png}
    \subcaption{TD-02}
  \end{minipage}
  \hspace{0.01\textwidth}
  \begin{minipage}[t]{0.30\textwidth}
    \centering
    \includegraphics[width=.8\textwidth]{assets/enshu_1/id3.png}
    \subcaption{TD-03}
  \end{minipage}
  \caption{演習1 テスト結果1}
  \label{fig:enshu1-result-1}
\end{figure}

\begin{figure}[H]
  \centering
  \begin{minipage}[t]{0.30\textwidth}
    \centering
    \includegraphics[width=.8\textwidth]{assets/enshu_1/id4.png}
    \subcaption{TD-04}
  \end{minipage}
  \hspace{0.01\textwidth}
  \begin{minipage}[t]{0.30\textwidth}
    \centering
    \includegraphics[width=.8\textwidth]{assets/enshu_1/id5.png}
    \subcaption{TD-05}
  \end{minipage}
  \hspace{0.01\textwidth}
  \begin{minipage}[t]{0.30\textwidth}
    \centering
    \includegraphics[width=.8\textwidth]{assets/enshu_1/id6.png}
    \subcaption{TD-06}
  \end{minipage}
  \caption{演習1 テスト結果2}
  \label{fig:enshu1-result-2}
\end{figure}

\begin{figure}[H]
  \centering
  \begin{minipage}[t]{0.30\textwidth}
  \end{minipage}
  \hspace{0.01\textwidth}
  \begin{minipage}[t]{0.30\textwidth}
    \centering
    \includegraphics[width=.8\textwidth]{assets/enshu_1/id7.png}
    \subcaption{TD-07}
  \end{minipage}
  \hspace{0.01\textwidth}
  \begin{minipage}[t]{0.30\textwidth}
  \end{minipage}
  \caption{演習1 テスト結果3}
  \label{fig:enshu1-result-3}
\end{figure}

画像より，TD-01 - 07までの結果は表\ref{tab:enshu1-tds}より，元から期待していたものに等しい．
そのため，仕様通りのソフトウェアを完成させられたと言える．

このテストは表\ref{tab:enshu1-decision}デシジョンテーブルに従って作成されたものだと言える．

\begin{table}[H]
  \centering
  \caption{演習1 デシジョンテーブル}
  \label{tab:enshu1-decision}
  \begin{tabular}{cccc}
    % ここにカラムを入力
    値1 & 値2 & 計算方法 & 出力 \\
    \hline
    % ここに各セルのデータを入力
    数値 & 数値 & 足し算 & 数値 \\
    数値 & 数値 & 引き算 & 数値 \\
    数値 & 数値 & かけ算 & 数値 \\
    数値 & 数値 & 割り算 & 数値 \\
    文字 & x & x & エラー \\
    x & 文字 & x & エラー \\
    数値 & 数値 & [空白] & エラー \\
  \end{tabular}
\end{table}

この時，xはドントケアである．

このテーブルは，四則演算プログラムの要件より抽出したものである．

値1，2自体は，有効クラスとして，数値(1, 0.002など)を持ち，無効クラスとして文字を持つ(aなど)．
また，計算方法は有効クラスとして4種の計算方法があり，無効クラスとして空白のみを持つ．
そのため，ただし，計算方法の有効クラスは全ての値がそれぞれ同値ではないため，個別にテストする必要がある．

以上のクラスタリング結果と，要件を踏まえテストケースを考えたい．

まず，値1, 2, 計算方法のすべてが有効クラスであれば，数値を出力を確認し(正常動作)，
入力のいずれかに無効クラスが含まれていれば，エラーを出力する(不正入力に対する正常動作)を確認することで，要件に必要となる動き全てを網羅したことになる．

従って，同値クラスの考え方を用いて，上記の表\ref{tab:enshu1-decision}を満たすテストケース(最低7個)を実行すれば，プログラムの動作が意図したものであるかどうかを判定することができるのである．

実際に，実行結果よりテストはすべて合格であったため，本四則演算プログラムは正常に動作していると，このテストから判断することができる．


\subsection{演習2 マルバツゲームのロジック実装，および拡張}
実際に表\ref{tab:enshu2-tds}に従い，テストケースを実行した結果が，図\ref{fig:enshu2-result-1}, \ref{fig:enshu2-result-2}である．

\begin{figure}[H]
  \centering
  \begin{minipage}[t]{0.45\textwidth}
    \centering
    \includegraphics[width=.8\textwidth]{assets/enshu_2/id1.png}
    \subcaption{TD-01}
  \end{minipage}\hspace{0.05\textwidth}\begin{minipage}[t]{0.45\textwidth}
    \centering
    \includegraphics[width=.8\textwidth]{assets/enshu_2/id2.png}
    \subcaption{TD-02}
  \end{minipage}
  \caption{演習2 テスト結果1}
  \label{fig:enshu2-result-1}
\end{figure}

\begin{figure}[H]
  \centering
  \begin{minipage}[t]{0.45\textwidth}
    \centering
    \includegraphics[width=.8\textwidth]{assets/enshu_2/id3.png}
    \subcaption{TD-03}
  \end{minipage}\hspace{0.05\textwidth}\begin{minipage}[t]{0.45\textwidth}
    \centering
    \includegraphics[width=.8\textwidth]{assets/enshu_2/id4.png}
    \subcaption{TD-04}
  \end{minipage}
  \caption{演習2 テスト結果2}
  \label{fig:enshu2-result-2}
\end{figure}

画像より，TD-01 - 04までの実行結果は，表\ref{tab:enshu2-tds}より意図したものであった．
そのため，仕様通りのソフトウェアを完成させたと言える．

次に，同値クラスの考え方を前提として，なぜこのテストでマルバツゲームの動作を検証できるのかを考える．

まず，盤面状態は2つの同値クラスに分けることができ，進行中の盤面と，全て埋まっている終了済みの盤面である．
また，ゲーム結果は3つの同値クラスに分けることができ，勝利，敗北，引き分けである．
さらに，プレイヤーは2つの同値クラスに分けることができ，先行と後行である．

そのため，単純に考えて$2\times 3\times 2$の12組テストが必要となる．
しかし，勝敗結果に関してはプレイヤーの片方を検証できれば，もう片方はその反対の結果になるはずである．
そのため，プレイヤーの状態を考慮しないものとすると，6組のテストケースのみを検証すれば良い．
しかし，進行中の盤面において，引き分けが発生することは論理的に考えられない．

加えて，三目並べは先行のプレイヤーが最後において負けることは発生しない．
なぜなら，8手目で後行が勝利しない場合，9手目で後行は行動できず，後行は負けか引き分け，先行は勝ちか引き分けにしかならないためである．

従って，同値クラスの考え方に基づき，以下の表\ref{tab:enshu2-cross}のテストケースのみを検証すれば良い．

\begin{table}[H]
  \centering
  \caption{マルバツゲーム テストケース}
  \label{tab:enshu2-cross}
  \begin{tabular}{c|ccc}
    & プレイヤーの勝利 & プレイヤーの敗北 & 引き分け \\
    \hline 
    進行中 & TD-01 & TD-02 & [発生しない] \\
    終了後 & TD-03 & [発生しない] & TD-04 \\
  \end{tabular}
\end{table}

本来は，発生しないと見込まれるテストケースに関しても，内部状態のダンプなどを用いて検証するべきであるが，ユーザーの入力が必ず先行になるという条件と，進行中の盤面において，論理的に引き分けだと言える盤面を制作できないという制約を鑑みて，
より，今回は無視して良いものとした．

表\ref{tab:enshu2-cross}に基づき制作された表\ref{tab:enshu2-tds}の結果は全て成功であったと言えるため，本プログラムの動作は，仕様書の意図した通りであると言える．


また，マルバツゲームの改良として市川先生の実験で作成したマルバツゲーム専用AIエージェントを実装し，CPUが選ぶ手の判断手段として，AIを用いる改良を施した．

実際の動作画面が図\ref{fig:enshu2-add}である．

\begin{figure}[H]
  \begin{center}
    \includegraphics[width=.5\textwidth]{assets/enshu_2_add/sample_1.png}
    \caption{改良版マルバツゲーム}
    \label{fig:enshu2-add}
  \end{center}  
\end{figure}

この時，各マスに挿入されている値は，最終的な勝ちやすさを示しており，大きいほど勝ちやすく，小さいほど負け，またはあいこになりやすくなる．

この値は，マルバツゲームを模した環境で，特定の報酬を用いてQ学習を行なったAIの指標である．

実際に，元のAIよりも勝率が上がり，あいこになるようになった．

\subsection{演習3 オセロゲームの実装，およびオセロゲームのテスト}
実際に表\ref{tab:enshu3-tds}に従ってテストを実行した結果が次の図\ref{fig:enshu3-result-1}, \ref{fig:enshu3-result-2}, \ref{fig:enshu3-result-3}である．

\begin{figure}[H]
  \centering
  \begin{minipage}[t]{0.45\textwidth}
    \centering
    \includegraphics[width=.98\textwidth]{assets/enshu_3/id1_s.png}
    \subcaption{TD-01 初期盤面}
  \end{minipage}\hspace{0.05\textwidth}\begin{minipage}[t]{0.45\textwidth}
    \centering
    \includegraphics[width=.98\textwidth]{assets/enshu_3/id1_e.png}
    \subcaption{TD-1 終了盤面}
  \end{minipage}
  \caption{TD-01 テスト結果}
  \label{fig:enshu3-result-1}
\end{figure}

\begin{figure}[H]
  \centering
  \begin{minipage}[t]{0.45\textwidth}
    \centering
    \includegraphics[width=.98\textwidth]{assets/enshu_3/id2_s.png}
    \subcaption{TD-02 初期盤面}
  \end{minipage}\hspace{0.05\textwidth}\begin{minipage}[t]{0.45\textwidth}
    \centering
    \includegraphics[width=.98\textwidth]{assets/enshu_3/id2_e.png}
    \subcaption{TD-02 終了盤面}
  \end{minipage}
  \caption{TD-02 テスト結果}
  \label{fig:enshu3-result-2}
\end{figure}

\begin{figure}[H]
  \centering
  \begin{minipage}[t]{0.45\textwidth}
    \centering
    \includegraphics[width=.98\textwidth]{assets/enshu_3/id3_s.png}
    \subcaption{TD-03 初期盤面}
  \end{minipage}\hspace{0.05\textwidth}\begin{minipage}[t]{0.45\textwidth}
    \centering
    \includegraphics[width=.98\textwidth]{assets/enshu_3/id3_e.png}
    \subcaption{TD-03 終了盤面}
  \end{minipage}
  \caption{TD-03 テスト結果}
  \label{fig:enshu3-result-3}
\end{figure}

まず，TD-01について考える．
このテストケースでは，初期配置として正常な初期配置(互い違いに中心4マスを黒白石で埋める)を用いて，一般的なオセロゲームのシナリオに従って，ゲームを進めることで終局するパターンを検証している．

一般的なオセロのルールに従って，互い違いに白黒石を置いていくことで，すべてのマス目が埋まった．
このとき，終局時の石の数と勝敗を正しく左上に表示できているため，マス目を残さず終局した場合の終了条件を正しく判定しているといえる．
したがって，テストケースに対してオセロゲームは正常に動作していると言える．

次に，TD-02について考える．
このテストケースでは，初期配置を意図的に設定し，盤面が白石のみになり，黒石の置ける場所がなくなることで，ゲームが終局するパターンを検証している．

初期状態のC2に黒石を置き，次のターンでAIがA2に白石が置かれることで，黒石の置く場所がなくなる．
その状態で，黒石のターンがパスされ，白石がD2に置かれ，すべての石が白になることで終局する．
結果として，白石のみで，あらなた白黒石の置ける場所がなくなるため，終局する．
このとき，左上に終局時の石の数と勝敗を正しく左上に表示できているため，マス目を残した状態で白石のみ残すという場合の終了条件を正しく判定していると言える．
したがって，テストケースに対してオセロゲームは正常に動作していると言える．

最後にTD-03について考える．
このテストケースでは，初期配置を意図的に設定し，盤面が黒石のみになり，白石の置ける場所がなくなることで，ゲームが終局するパターンを検証している．

初期状態のH1に黒石を置くことで，すべての石が黒色になる．
結果として，黒石のみで，あらなた白黒石の置ける場所がなくなるため，終局する．
このとき，左上に終局時の石の数と勝敗を正しく左上に表示できているため，マス目を残した状態で黒石のみ残すという場合の終了条件を正しく判定していると言える．
したがって，テストケースに対してオセロゲームは正常に動作していると言える．

\subsection{演習4 ユニットテストのテストケースの制作方法}
テスト対象のメソッドの入出力要件は，表\ref{tab:enshu4-target-method-1}, \ref{tab:enshu4-target-method-2}に従う．
この時，同値クラスは次の表\ref{tab:enshu4-class-1}, \ref{tab:enshu4-class-2}のようになる．

\begin{table}[H]
  \centering
  \caption{演習4 同値クラス(点数)}
  \label{tab:enshu4-class-1}
  \begin{tabular}{ccc}
    % ここにカラムを入力
    数値 & 有効/無効 & クラス \\
    \hline
    % ここに各セルのデータを入力
    $101 \leq x $ & 無効 & \\
    $80 \leq x \leq 100$ & 有効 & 優 \\
    $65 \leq x \leq 79$ & 有効 & 良 \\
    $60 \leq x \leq 64$ & 有効 & 可 \\
    $40 \leq x \leq 59$ & 有効 & 不可A \\
    $0 \leq x \leq 39$ & 有効 & 不可B \\
    $x \leq 0$ & 無効 & \\
  \end{tabular}
\end{table}

\begin{table}[H]
  \centering
  \caption{演習4 同値クラス(欠課)}
  \label{tab:enshu4-class-2}
  \begin{tabular}{ccc}
    % ここにカラムを入力
    数値 & 有効/無効 & クラス \\
    \hline
    % ここに各セルのデータを入力
    $31 \leq y $ & 無効 & \\
    $15 \leq y \leq 30$ & 有効 & x(バツ) \\
    $0 \leq y \leq 14$ & 有効 & 点数の出力に依存 \\
    $y \leq -1$ & 無効 &  \\
  \end{tabular}
\end{table}

この時，点数入力の同値クラスは7クラス存在し，欠課の同値クラスは4クラス存在する．
従って，同値分割法により，各同値クラスから代表値を1つ取り出してテストした場合，テストケースの数$C_1$は次のようになる．

\begin{equation}
  C_1 = 1 \times 7 \times 4 = 28
\end{equation}

また，そのため境界値は点数入力において12個，欠課において6個存在する．
従って，境界値分析を行う場合，必要なテストケースの数$C_2$は次のようになる．

\begin{equation}
  C_2 = 12 \times 6 = 72
\end{equation}

テスト対象のメソッドでは，値の範囲において出力が大きく変化する．
そのため，値の範囲に着目した境界値テストを採用する．

従って，必要なテストケースの数は$C_2 = 72$より72個である．



\subsection{演習5 自前の実装プログラムのテスト}
\subsubsection{Cellクラス コンストラクタ}
まず，テストケースの作成を行なった．
本コンストラクタは，RowとColumnの入力値を持ち，それぞれ正常な値の範囲が$0 \leq Row, Column \leq 7$である．
また，これらの値は設計書\cite{journal:sekkeisyo}よりint型で宣言されていることが明らかであるため，実際には整数の上記範囲が正常な値の範囲である．

従ってRowにおける同値分割法による同値クラスは以下の表\ref{tab:enshu5-cell-const-1}のようになる．

\begin{table}[H]
  \centering
  \caption{Rowパラメーターの同値クラス}
  \label{tab:enshu5-cell-const-1}
  \begin{tabular}{cc}
    % ここにカラムを入力
    範囲 & 有効/無効 \\
    \hline
    % ここに各セルのデータを入力
    $Row \leq -1$ & 無効 \\
    $0 \leq Row \leq 7$ & 有効 \\
    $8 \leq Row$ & 無効 \\
  \end{tabular}
\end{table}

また，Columnにおいて，同値分割法による同値クラスは以下の表\ref{tab:enshu5-cell-const-2}のようになる．

\begin{table}[H]
  \centering
  \caption{Columnパメーターの同値クラス}
  \label{tab:enshu5-cell-const-2}
  \begin{tabular}{cc}
    % ここにカラムを入力
    範囲 & 有効/無効 \\
    \hline
    $Column \leq -1$ & 無効 \\
    $0 \leq Column \leq 7$ & 有効 \\
    $8 \leq Column$ & 無効 \\
  \end{tabular}
\end{table}

本コンストラクタの正常な入力に対する正常動作と異常入力に対する正常動作はを考えたい．

これは，以下の設計書\cite{journal:sekkeisyo}より導く．

\begin{quote}
public Cell(int row, int column) \\

コンストラクタ．Row プロパティと Column プロパティを引数で初期化する．他のプロパティは，ピースが
置かれておらず，置けない状態を表すように初期化する．row や column のいずれかが範囲を超える場合，Row
と Column の両方を-1 にする．
\end{quote}

特に，ピースが置かれていない状態はや置けない状態の表し方に関しては，
本クラスのメンバ変数にIsBlack, IsBlank, CanBlackPlace, CanWhitePlaceが存在する．

各メンバの意味は次の表\ref{tab:cell-member-variable}である．

\begin{table}[H]
  \centering
  \caption{Cellクラスのメンバ}
  \label{tab:cell-member-variable}
  \begin{tabular}{ccc}
    % ここにカラムを入力
    型 & メンバ名 & 摘要 \\
    \hline
    % ここに各セルのデータを入力
    public bool & IsBlank & 空白か．空白だとTrue, 埋まっていればFalse \\
    public bool & IsBlack & 黒色か．白色または空白だとFalse，黒だとTrue \\
    public bool & CanBlackPlace & 黒色を置けるか．置けるとTrue，置けるならFalse \\
    public bool & CanWhitePlace & 白色を置けるか．置けるとTrue，置けるならFalse \\
  \end{tabular}
\end{table}

以上を踏まえ，正常な入力値と異常な入力値に対して取られる正常な動作は，次の表\ref{tab:enshu5-cell-const}である．

\begin{table}[H]
  \centering
  \caption{コンストラクタの動作表}
  \label{tab:enshu5-cell-const}
  \begin{tabular}{ccccccc}
    % ここにカラムを入力
    入力値 & Row & Column & IsBlank & IsBlack & CanBalckPlace & CanWhitePlace\\
    \hline
    % ここに各セルのデータを入力
    正常 & 入力値 & 入力値 & False & False & False & False\\
    異常 & -1 & -1 & False & False & False & False\\
  \end{tabular}
\end{table}

このとき，このコンストラクタの動作として，境界値付近の動作の性質が変化点が重要である．

従って，RowとColumnに対してそれぞれの境界値をクロスさせたテストケースを網羅的に実行し，先の動作表より推測される動作と一致しているかを検証する
テストを行うことで，動作の正常性を示すこととした．

この際，RowとColumnの境界値をクロスさせるために必要なテストケースおよび，その実行結果は次の表\ref{tab:enshu5-cell-const-tds}となった．

この表における合格とは，期待される動作に一致したことを示す．

\begin{table}[H]
  \centering
  \caption{Cellコンストラクタのテストケース}
  \label{tab:enshu5-cell-const-tds}
  \begin{tabular}{ccccc}
    % ここにカラムを入力
    ID & Row & Column & 期待される動作 & 結果 \\
    \hline
    % ここに各セルのデータを入力
    TD-01 & 0  & 0  & 正常 & 合格 \\
    TD-02 & 0  & 7  & 正常 & 合格 \\
    TD-03 & 0  & -1 & 異常 & 合格 \\
    TD-04 & 0  & 8  & 異常 & 合格 \\
    TD-05 & 7  & 0  & 正常 & 合格 \\
    TD-06 & 7  & 7  & 正常 & 合格 \\
    TD-07 & 7  & -1 & 異常 & 合格 \\
    TD-08 & 7  & 8  & 異常 & 合格 \\
    TD-09 & -1 & 0  & 異常 & 合格 \\
    TD-10 & -1 & 7  & 異常 & 合格 \\
    TD-11 & -1 & -1 & 異常 & 合格 \\
    TD-12 & -1 & 8  & 異常 & 合格 \\
    TD-13 & 8  & 0  & 異常 & 合格 \\
    TD-14 & 8  & 7  & 異常 & 合格 \\
    TD-15 & 8  & -1 & 異常 & 合格 \\
    TD-16 & 8  & 8  & 異常 & 合格 \\
  \end{tabular}
\end{table}

実際にこれらのデータをMSTestライブラリを用いて，テストメソッドを作成し，テストした．

最終的に全テストケースが合格したため，本テストは成功であり，プログラミングしたコンストラクタは正常な動作を行なっていたと言える．



\subsubsection{Cellクラス Setメソッド}
まず，テストケースの作成を行なった．

Setメソッドの入力および動作要件は以下のとおりである．

\begin{quote}
public void Set(bool isBlack)\\

セルにピースを置く．反転による色の変化もこのメソッドで実現するため，既に置かれているセルに対し
ても処理を行う（CanBlackPlace や CanWhitePlace の値はこのメソッドの処理に影響しない）．
\end{quote}

本メソッドでは，セルにピースを置く，または反転したセルの状態を記録するという2つの役割を行う．
つまり，どんな状態であってもisBlackの色をセルにセットするという役割を担う．

このとき，入力値はメソッド自体のisBlackおよび，CellクラスのメンバであるIsBlack, IsBlankの3種類である．
なぜなら，どんな状態のセルにもピースを置くのが重要なのであって，RowやColumn設定された値を無視することができるからである．
このメソッドでは，CanBlackPlace, CanWhitePlaceの値はこのメソッドの処理に影響しないことが明記されているため，入力値から除外している．

従って，$2\times 2\times 2=8$ケーステストすればすべての値を網羅できることとなる．

そのため，次の表\ref{tab:enshu5-cell-set-tds}のような入力値網羅のテストケースを作成した．

この表は横軸が(IsBlack, IsBlank)であり，縦軸がisBlackである．

\begin{table}[H]
  \centering
  \caption{Cellクラス Setメソッド テストケース}
  \label{tab:enshu5-cell-set-tds}
  \begin{tabular}{c|cccc}
    % ここにカラムを入力
    & (False, False) & (True, False) & (True, True), & (False, True) \\
    \hline
    % ここに各セルのデータを入力
    False & TD-01 & TD-02 & - & TD-03 \\
    True & TD-04 & TD-05 & - & TD-06 \\
  \end{tabular}
\end{table}

このとき，IsblankがTrueでIsBlackがTrueな状況がありえない点に関してだが，これはIsBlankがTrueであるのは何もおいていない状態を示しており，
その上で，何もにおいていない時はIsBlackもFalseになるという仕様がCellのコンストラクタに存在し，それ以外の点でIsBlankがTrueとなることはないためである．
また，その動作は先のコンストラクタのテストで保証されている．

当然，外的要因によって上書きされた場合は，その通りではないものの，仕様上に存在しない値に対しては，テストしないものとする．

各テストケースTD-01-06までの期待動作は次の表\ref{tab:enshu5-cell-set-exp}の通りである．

\begin{table}[H]
  \centering
  \caption{テストケース TD-01-06までの期待動作}
  \label{tab:enshu5-cell-set-exp}
  \begin{tabular}{ccc}
    % ここにカラムを入力
    ID & IsBlack & IsBlank \\
    \hline
    % ここに各セルのデータを入力
    TD-01 & False & False \\
    TD-02 & False & False \\
    TD-03 & False & False \\
    TD-04 & True & False \\
    TD-05 & True & False \\
    TD-06 & True & False \\
  \end{tabular}
\end{table}

この時，反転による色の変化をこのメソッドで実行するというのが非常に重要である．
つまり，このメソッドはCellインスタンスの状況に応じて動作が変わるため，ステートマシンとしての側面を持つことになる．
そのため，その側面をテストする．

具体的には次の状態遷移が考えられる．
状態は(IsBlack, IsBlank)の組みである．

\begin{enumerate}
  \item (False, True)
  \item (False, False)
  \item (True, False)
  \item (True, False)
  \item 以下(False, False), (True, False)を任意の確率で繰り返し状態遷移
\end{enumerate}

\begin{enumerate}
  \item (False, True)
  \item (True, False)
  \item (False, False)
  \item (False, False)
  \item 以下(False, False), (True, False)を任意の確率で繰り返し状態遷移
\end{enumerate}

本来であれば，任意の割合で状態遷移する部分をすべてテストする必要がある．
しかし，この2組の状態遷移をテストすべてテストすれば，初期状態以降の状態遷移は実質的な繰り返しであり，すでにテストされた結果を適用できるため，
テスト不要になる．

従って，本テストではIsBlankとIsBlackという内部状態にisBlackという外的な3種類の入力値に対して，状態遷移的に出力値を検証することで，
すべての動作パターンに対してテストを行なうことに等しいと判断し，メソッドの動作を検証する，というテストケースを作成する．

そのため，今回は2つのテストデータ(2種類の状態遷移の確認)を用意し，それらを連続的に処理し，その際の値の変化を見る．

具体的に，次のようになる．

\begin{enumerate}
  \item 初期状態のCell Aを用意
  \item Aに黒をセット
  \item Aの内部状態を確認 (TD-06に相当)
  \item Aに白をセット
  \item Aの内部状態を確認 (TD-02に相当)
  \item Aに白をセット
  \item Aの内部状態を確認 (TD-01に相当)
\end{enumerate}

\begin{enumerate}
  \item 初期状態のCell Bを用意
  \item Bに白をセット
  \item Bの内部状態を確認 (TD-03に相当)
  \item Bに黒をセット
  \item Bの内部状態を確認 (TD-04に相当)
  \item Bに黒をセット
  \item Bの内部状態を確認 (TD-05に相当)
\end{enumerate}

実際に，以上の内容を用いてテストケースを作成し，TD-01からTD-06までの出力をすべて確認した．
結果，それらはすべて合格であったため，本メソッドの動作は仕様書に照らして正しいと言える．


\subsubsection{Boardクラス CalcConditionメソッド テストケース作成}
まず，テストケースを考えた．

このメソッドの入力はBoardにあるCellの全状態である．
また，動作として，Board上の64個のセルに対するCanWhitePlace, CanBlackPlaceの変更および，Board自体のCanWhitePlace，CanBlackPlaceの変更である．

今回は，特定のCellの状態を同値クラスとみなし，テスト対象のCellを減らしつつ，CanWhitePlaceとCanBlackPlaceを完全に独立した指標としてとらえ，テストケースを設計していきたい．

また，今回使用するメソッドの仕様書は次の通りである．

\begin{quote}
public void CalcCondition() \\

ボード上の全てのセルについて，黒と白それぞれがそのセルにピースを置けるかどうか計算し，Cells に格
納された各セルの情報を更新する．各マスにおけるかどうかは，1)そのマスが空いているか，2) 8 方向のい
ずれかが挟めるかを確認し，判断する．「挟めるか」の判定は CanFlip メソッドを用いる．また，Board クラ
スのプロパティである CanBlackPlace，CanWhitePlace もここで更新する．
\end{quote}

以下，特筆すべき事項がない限り，本文言から入力，動作要件を導出しているものとする．

まず，BoardのCanWhitePlace/CanBlackPlaceへのテスト計画を考える．
可読性を上げるため，今回はCanWhitePlaceに関してのみテストケースの概略を記述するが，CanBlackPlaceに関しても同様の手順でテストケースを設計することができる．

まず，CanWhitePlaceはBoard上で一箇所以上白石を置ける場所が存在するかを表している．
つまり盤面の種別に関係なく，置ける場所が0箇所であれば，CanWhitePlaceはFalseになり，一箇所以上あればTrueに変化する．

よって以下のようにBoardの状態を同値クラスに分ける．

\begin{itemize}
  \item 白石を一箇所もおけない
  \item 白石を一箇所以上置ける
\end{itemize}

同様の考え方で，CanBlackPlaceに関しても，以下のようにBoardの状態を同値クラスに分ける．

\begin{itemize}
  \item 黒石を一箇所もおけない
  \item 黒石を一箇所以上置ける
\end{itemize}

従って，CanWhitePlaceを検証するのに2つのテストケースが必要となり，
かつCanBlackPlaceを検証するのに2つのテストケースが必要となる．

従って，CanWhitePlaceとCanBlackPlaceのデシジョンテーブルは次の表\ref{tab:enshu5-board-decision-board}のようになる．

\begin{table}[H]
  \centering
  \caption{BoardのCanWhite/BlackPlaceのデシジョンテーブル}
  \label{tab:enshu5-board-decision-board}
  \begin{tabular}{cccc}
    % ここにカラムを入力
    ID & 盤面状態 & CanBlackPlace & CanWhitePlace \\ 
    \hline
    % ここに各セルのデータを入力
    37 & 黒石が一箇所以上置ける & True & x \\
    38 & 黒石が一箇所もおけない & False & x \\
    39 & 白石が一箇所以上置ける & x & True \\
    40 & 白石が一箇所もおけない & x & False \\
  \end{tabular}
\end{table}

このとき，xはドントケアである．



次に，各CellのCanWhitePlace/CanBlackPlaceに関してへのテスト計画を考える．
同様に，白石，CanWhitePlaceについてのみ概略を記述するものの，黒石，CanBlackPlaceに関しても同様にテストケースを設計できる．

まず，すべての盤面においてCellを確認することが理想である．
しかし，この理想ではテスト項目が$3^{64}$となってしまい，あまりに膨大である．

そのため，テスト対象となる盤面状態からCellの特徴に着目し，同値分割を行い，同値分割法によりテストを行う．

概略として，あるCellの状態に対する動作が正しいと言えた場合，類似する状態にあるCell全ての動作を正しいということができ，
それをすべての盤面のすべてのCellに対して個別に適用することで，すべての盤面に対するCalcConditionの動作が正しいということを目指す．

このあるCellの状態をいずれかの手法でクラスタリングし，同値クラスに分割する．
最終的に，各同値クラスから代表値を算出し，テストを完了する．

まず，CanWhitePlaceは次の要件に従って変更される．

\begin{itemize}
  \item 検査対象のマスの空白状態(IsBlank)
  \item 8方向のいずれかで挟める状態(CanFlip)
\end{itemize}

この2つの要件に従って，CanWhitePlaceは変更される．
この時，検査対象マスのCanWhitePlaceに関するデシジョンテーブルは次の表\ref{tab:enshu5-board-decision-white}のように定義される．

\begin{table}[H]
  \centering
  \caption{CanWhitePlaceの変更要件}
  \label{tab:enshu5-board-decision-white}
  \begin{tabular}{cccc}
    % ここにカラムを入力
    ID & 空白 & 白に関して8方向いずれかで挟める & CanWhitePlace \\
    \hline
    % ここに各セルのデータを入力
    19 & 空白 & 真右で挟める & True \\
    20 & 空白 & 真下で挟める & True \\
    21 & 空白 & 真左で挟める & True \\
    22 & 空白 & 真上で挟める & True \\
    23 & 空白 & 右下で挟める & True \\
    24 & 空白 & 左下で挟める & True \\
    25 & 空白 & 左上で挟める & True \\
    26 & 空白 & 右上で挟める & True \\
    27 & 空白 & いずれも挟めない & False \\
    28 & うまっている & 真右で挟める & False \\
    29 & うまっている & 真下で挟める & False \\
    30 & うまっている & 真左で挟める & False \\
    31 & うまっている & 真上で挟める & False \\
    32 & うまっている & 右下で挟める & False \\
    33 & うまっている & 左下で挟める & False \\
    34 & うまっている & 左上で挟める & False \\
    35 & うまっている & 右上で挟める & False \\
    36 & うまっている & いずれも挟めない & False \\
  \end{tabular}
\end{table}

この表より，すべてのCellの状態は以上の18個の同値クラスに分割することができるとわかった．

そのため，各デシジョンテーブルと合致するテストケースを作成し，そのテストの結果を参照することで，CanWhitePlaceの出力が正しいのか否かを判定することができる．

また，同様の考え方を黒石，CanBlackPlaceにも適用する．
すると，検査対象マスのCanBlackPlaceに関するデシジョンテーブルは次の表\ref{tab:enshu5-board-decision-black}のように定義される．

\begin{table}[H]
  \centering
  \caption{CanBlackPlaceの変更要件}
  \label{tab:enshu5-board-decision-black}
  \begin{tabular}{cccc}
    % ここにカラムを入力
    ID & 空白 & 黒に関して8方向いずれかで挟める & CanBlackPlace \\
    \hline
    % ここに各セルのデータを入力
    01 & 空白 & 真右で挟める & True \\
    02 & 空白 & 真下で挟める & True \\
    03 & 空白 & 真左で挟める & True \\
    04 & 空白 & 真上で挟める & True \\
    05 & 空白 & 右下で挟める & True \\
    06 & 空白 & 左下で挟める & True \\
    07 & 空白 & 左上で挟める & True \\
    08 & 空白 & 右上で挟める & True \\
    09 & 空白 & いずれも挟めない & False \\
    10 & うまっている & 真右で挟める & False \\
    11 & うまっている & 真下で挟める & False \\
    12 & うまっている & 真左で挟める & False \\
    13 & うまっている & 真上で挟める & False \\
    14 & うまっている & 右下で挟める & False \\
    15 & うまっている & 左下で挟める & False \\
    16 & うまっている & 左上で挟める & False \\
    17 & うまっている & 右上で挟める & False \\
    18 & うまっている & いずれも挟めない & False \\
  \end{tabular}
\end{table}


以上の表\ref{tab:enshu5-board-decision-white}, \ref{tab:enshu5-board-decision-black}について検証することであるセルに関して，次のことがわかる．

\begin{itemize}
  \item 黒/白いずれかの状態で
  \item かつ，空白／すでに置かれている状態で
  \item かつ，8方向のある方向のみに挟める/いずれも挟めない
  \item という状態のセルに対するCalcConditionの動作が正しいこと
\end{itemize}

また，すべての盤面に存在するCellの状態は先述したID01からID36までのデシジョンテーブルの要件のいずれかに属するため，これらを同値クラスとして断定できる．
これはID01からID36が示すクラスが，あるセルが取りうる状態が先述した，2つの判定要件のみであり，その2つの要件をクロスして作成したクラスであるためである．


以上より，同値クラスの考え方に基づいて，CalcConditionのメソッドの正確性を占めるテストケースを考えることができた．

しかし，今回のCalcConditionは，盤面上すべての値に対して作用する必要があるが，同値クラスの考え方に基づいたテストケースでは，その点に関して考慮できていない．
そのため，実際のテストデータの作成では，境界値である端付近でテストとデータを収集し，端付近位おいても正しく動作しているのか否かを検証する．


\subsubsection{Boardクラス CalcConditionメソッド テストデータ作成・実行}
先の表\ref{tab:enshu5-board-decision-board}, \ref{tab:enshu5-board-decision-white}, \ref{tab:enshu5-board-decision-black}の結果を踏まえて，
これらのテストケースをすべて網羅可能であるかつ，ボードにおける境界値分析が行えるようなテストデータを作成した．

実際のテストデータは，次の図\ref{fig:enshu5-td-1}から\ref{fig:enshu5-td-6}である．


\begin{figure}[H]
  \centering
  \begin{minipage}[t]{0.30\textwidth}
    \centering
    \includegraphics[width=.8\textwidth]{assets/enshu_5/calc_condition/TD-01.png}
    \subcaption{TD-01}
  \end{minipage}
  \hspace{0.01\textwidth}
  \begin{minipage}[t]{0.30\textwidth}
    \centering
    \includegraphics[width=.8\textwidth]{assets/enshu_5/calc_condition/TD-02.png}
    \subcaption{TD-02}
  \end{minipage}
  \hspace{0.01\textwidth}
  \begin{minipage}[t]{0.30\textwidth}
    \centering
    \includegraphics[width=.8\textwidth]{assets/enshu_5/calc_condition/TD-03.png}
    \subcaption{TD-03}
  \end{minipage}
  \caption{演習5 テストデータ1}
  \label{fig:enshu5-td-1}
\end{figure}

\begin{figure}[H]
  \centering
  \begin{minipage}[t]{0.30\textwidth}
    \centering
    \includegraphics[width=.8\textwidth]{assets/enshu_5/calc_condition/TD-04.png}
    \subcaption{TD-04}
  \end{minipage}
  \hspace{0.01\textwidth}
  \begin{minipage}[t]{0.30\textwidth}
    \centering
    \includegraphics[width=.8\textwidth]{assets/enshu_5/calc_condition/TD-05.png}
    \subcaption{TD-05}
  \end{minipage}
  \hspace{0.01\textwidth}
  \begin{minipage}[t]{0.30\textwidth}
    \centering
    \includegraphics[width=.8\textwidth]{assets/enshu_5/calc_condition/TD-06.png}
    \subcaption{TD-06}
  \end{minipage}
  \caption{演習5 テストデータ2}
  \label{fig:enshu5-td-2}
\end{figure}

\begin{figure}[H]
  \centering
  \begin{minipage}[t]{0.30\textwidth}
    \centering
    \includegraphics[width=.8\textwidth]{assets/enshu_5/calc_condition/TD-07.png}
    \subcaption{TD-07}
  \end{minipage}
  \hspace{0.01\textwidth}
  \begin{minipage}[t]{0.30\textwidth}
    \centering
    \includegraphics[width=.8\textwidth]{assets/enshu_5/calc_condition/TD-08.png}
    \subcaption{TD-08}
  \end{minipage}
  \hspace{0.01\textwidth}
  \begin{minipage}[t]{0.30\textwidth}
    \centering
    \includegraphics[width=.8\textwidth]{assets/enshu_5/calc_condition/TD-09.png}
    \subcaption{TD-09}
  \end{minipage}
  \caption{演習5 テストデータ3}
  \label{fig:enshu5-td-3}
\end{figure}

\begin{figure}[H]
  \centering
  \begin{minipage}[t]{0.30\textwidth}
    \centering
    \includegraphics[width=.8\textwidth]{assets/enshu_5/calc_condition/TD-10.png}
    \subcaption{TD-10}
  \end{minipage}
  \hspace{0.01\textwidth}
  \begin{minipage}[t]{0.30\textwidth}
    \centering
    \includegraphics[width=.8\textwidth]{assets/enshu_5/calc_condition/TD-11.png}
    \subcaption{TD-11}
  \end{minipage}
  \hspace{0.01\textwidth}
  \begin{minipage}[t]{0.30\textwidth}
    \centering
    \includegraphics[width=.8\textwidth]{assets/enshu_5/calc_condition/TD-12.png}
    \subcaption{TD-12}
  \end{minipage}
  \caption{演習5 テストデータ4}
  \label{fig:enshu5-td-4}
\end{figure}

\begin{figure}[H]
  \centering
  \begin{minipage}[t]{0.30\textwidth}
    \centering
    \includegraphics[width=.8\textwidth]{assets/enshu_5/calc_condition/TD-13.png}
    \subcaption{TD-13}
  \end{minipage}
  \hspace{0.01\textwidth}
  \begin{minipage}[t]{0.30\textwidth}
    \centering
    \includegraphics[width=.8\textwidth]{assets/enshu_5/calc_condition/TD-14.png}
    \subcaption{TD-14}
  \end{minipage}
  \hspace{0.01\textwidth}
  \begin{minipage}[t]{0.30\textwidth}
    \centering
    \includegraphics[width=.8\textwidth]{assets/enshu_5/calc_condition/TD-15.png}
    \subcaption{TD-15}
  \end{minipage}
  \caption{演習5 テストデータ5}
  \label{fig:enshu5-td-5}
\end{figure}

\begin{figure}[H]
  \centering
  \begin{minipage}[t]{0.30\textwidth}
    \centering
    \includegraphics[width=.8\textwidth]{assets/enshu_5/calc_condition/TD-16.png}
    \subcaption{TD-16}
  \end{minipage}
  \caption{演習5 テストデータ6}
  \label{fig:enshu5-td-6}
\end{figure}

今回のテストデータ設計では，TD-01とTD-02というように2組のテストデータがセットでテストされることで意味のあるテストデータを作成している．
赤丸のマークはそのテストケースにおける検査対象のマスを表している．
詳細に関しては後述する．
また以降，盤面上の位置を，縦軸をRow, 横軸をColumnとし(Row, Column)の順番で記載し，左上を(0, 0), 右上を(0, 7), 右下を(7, 7), 左下を(7, 0)と記述する．

TDから始まる各テストデータと，表\ref{tab:enshu5-board-decision-board}, 
\ref{tab:enshu5-board-decision-white}, \ref{tab:enshu5-board-decision-black}の対応表が表\ref{tab:enshu5-board-tds}である．

\begin{table}[H]
  \centering
  \caption{テストデータとテストケースの網羅性}
  \label{tab:enshu5-board-tds}
  \begin{tabular}{cc}
    % ここにカラムを入力
    テストデータID & テストケースID \\
    \hline
    % ここに各セルのデータを入力
    TD-01 & 1, 2, 3, 4 \\
    TD-02 & 10, 11, 12, 13 \\
    TD-03 & 5, 6, 7, 8 \\
    TD-04 & 14, 15, 16, 17 \\
    TD-05 & 9 \\
    TD-06 & 18 \\
    TD-07 & 19, 20, 21, 22 \\
    TD-08 & 28, 29, 30, 31 \\
    TD-09 & 23, 24, 25, 26 \\
    TD-10 & 32, 33, 34, 35 \\
    TD-11 & 27 \\
    TD-12 & 36 \\
    TD-13 & 37 \\
    TD-14 & 38 \\
    TD-15 & 39 \\
    TD-16 & 40 \\
  \end{tabular}
\end{table}

各テストケースを説明する．

\paragraph{A: (TD-01, TD-02), B: (TD-07, TD-08)}

(TD-01, TD-02)の組みは黒に関するCellのCanBlackPlaceの性質のみをテストし，(TD-07, TD-08)の組みは白に関するCellのCanWhitePlaceの性質のみをテストする．
また，これらの組みは石の色に関して全くの反転した盤面であるため，(TD-01, TD-02) (以下，Aと呼ぶ)のみ概略を記載する．

Aの検査対象マスは，(0, 0), (0, 7), (7, 7), (7, 0)である．
検査対象マスを端に置いている理由として，CalcConditionの動作範囲はRow Columnともに0から7までであり，この数値が境界値となるためである．
境界値分析を用いて，Board上の境界値に対して動作が正しいかを同時に確認するために，検査対象マスを4隅に設定した．

まず，TD-01について考える．
このとき，(0, 0)のマスの状況としては，空白であり，真右方向にのみ挟める状況である．
つまり，(0, 0)のCanBlackPlaceがTrueとなる場合，この変化への作用素は真右方向の挟めるという状況のみである．
従って，このテストケースの実行結果が正しいとなると，テストケースID1の状況に対してCalcConditionがCellのCanBlackPlaceへ与える影響は正しいと言える．

同様の手順で(0, 7)は空白，真下にのみ挟める状況に，(7, 7)は空白，真左にのみ挟める状況に，(7, 0)は空白，真上にのみ挟める状況を表し，
それらの状況の動作が正しいか検証するために検査する
(テストケースIDはそれぞれ，2，3，4に対応)．

反対に，TD-02について考える．
このとき，(0, 0)のマスの状況としては，埋まっており，真右方向にのみ挟める状況である．
また，TD-01とTD-02の(0, 0)におけるマスの状況の差は，空白か埋まっているかのみであり，その他の状況は全く同じである．

つまり，(0, 0)のCanBlackPlaceはTD-01の状況ではTrueとなっているが，TD-02でFalseに変化した場合，
この変化への作用素は埋まっているという状況のみとなる．
従って，このテストケースの実行結果が正しいとなると，テストケースID10の状況に対してCalcConditionがCellのCanBlackPlaceへ与える影響は正しいと言える．

他の検査対象マスに関しても同様である(テストケースIDはそれぞれ，11，12，13に対応)．


従って，Aのすべてのテストケースが正しいとすれば，あるCellがいずれかの空白状況であり，
真右，真下，真左，真上のいずれかで挟める状況に対するCellのCanBlackPlaceへの作用が正しいと言える．

Bに関しては，これが全く反対になり，上記と同様の条件のCellに対してCalcConditionがCellのCanWhitePlaceへの作用が正しいと言える
(テストケースIDはそれぞれ，19, 20, 21, 22，28, 29, 30, 31に対応)．


\paragraph{A: (TD-03, TD-05), B: (TD-09, TD-10)}
先ほどと同様にAのケースについてのみ考え，Bはその反転であるため記載しない．

Aの検査対象マスは先ほどと同様の理由で，(0, 0), (0, 7), (7, 7), (7, 0)である．

まず，TD-03を考えたい．
このとき，(0, 0)のマスの状況としては，空白であり，右斜め下方向にのみ挟める状況である．
つまり，(0, 0)のCanBlackPlaceがTrueとなれば，この変化への作用素は右斜め下方向に挟める状況のみとなる．
従って，このテストケースの実行結果が正しいとなる，テストケースID5の状況に対して，CalcConditionがCellのCanBlackPlaceへ与える影響は正しいと言える．

同様の手順で(0, 7)は空白，左下にのみ挟める状況に, (7, 7)は空白，左上のみ挟める状況に, (7, 0)は空白，右上のみ挟める状況を表し，
それらの状況の動作が正しいか検証するために検査する
(テストケースIDはそれぞれ6, 7, 8に対応する)．

反対に，TD-04について考える．
このとき，(0, 0)のマスの状況としては，埋まっており，右斜め下方向にのみ挟める状況である．
また，TD-03とTD-04の(0, 0)におけるマスの状況の差は，空白か埋まっているかのみであり，その他の状況は全く同じである．

つまり，(0, 0)のCanBlackPlaceはTD-03の状況ではTrueとなっているが，TD-04でFalseに変化した場合，
この変化への作用素は埋まっているという状況のみとなる．
従って，このテストケースの実行結果が正しいとなると，テストケースID14の状況に対してCalcConditionがCellのCanBlackPlaceへ与える影響は正しいと言える．

他の検査対象マスに関しても同様である(テストケースIDはそれぞれ，14，15，16，17に対応)．

従って，Aのすべてのテストケースが正しいとすれば，あるCellがいずれかの空白の状況であり，
右斜め下，左斜め下，左斜め上，右斜め上のいずれかで挟める状況に対するCellのCanBlackPlaceへの作用が正しいと言える．

Bに関しては，これが全く反対になり，上記と同様の条件のCellに対してCalcConditionがCellのCanWhitePlaceへの作用が正しいと言える
(テストケースIDはそれぞれ，23, 24, 25, 26，32, 33, 34, 35に対応)．


\paragraph{A: (TD-05, TD-06), B: (TD-11, TD-12)}
先ほどと同様に，Aのケースについてのみ考え，Bはその反転であるため記載しない．

Aの検査対象マスは先ほどと同様の理由で，(0, 0), (0, 7), (7, 7), (7, 0)である．

まず，TD-05を考えたい．
これは，すべての検査対象マスが，空白で，いずれの方向に対しても挟めない状況にある．
このとき，すべての検査対象マスのCanBlackPlaceはFalseとなる．
なぜなら，挟める場所がないからである．

従って，このテストケースの実行結果が正しいとなると，テストケースID9の状況に対してCalcConditionがCellのCanBlackPlaceへ与える影響は正しいと言える．

次に，TD-06を考えたい．
これは，すべての検査対象マスが，埋まっており，いずれの方向に対しても挟めない状況にある．
また，TD-05とTD-06の差はマスの空白状況のみである．

つまり，TD-06の検査対象マスのCanBlackPlaceがFalseであれば，埋まっているという状況によってFalseになったと言える．

従って，このテストケースの実行結果が正しいとなると，テストケースID18の状況に対してCalcConditionがCellのCanBlackPlaceへ与える影響は正しいと言える．

Bに関しては，これが全く反対になり，上記と同様の条件のCellに対してCalcConditionがCellのCanWhitePlaceへの作用が正しいと言える
(テストケースIDはそれぞれ，27, 36に対応)．

\paragraph{A: (TD-13, TD-14), B: (TD-15, TD-16)}
先ほどと同様に，Aのケースについてのみ考え，Bはその反転であるため記載しない．

Aの検査対象はBoardのCanBlackPlaceである．

まず，TD-13について考える．
この時，盤面上には赤丸の部分にだけ黒石を置くことができるため，置ける場所は1である．
つまり，CanBlackPlaceはTrueとなる．

従って，このテストケースの実行結果が正しいとなると，テストケースID37の状況(境界値分析)において，CalcConditionがBoardのCanBlackPlaceへ与える影響は正しいと言える．

次に，TD-14を考えたい．
この時，盤面上には黒石を一箇所も置くことができない．
また，TD-13とTD-14の差は，TD-13で置くことができた黒石の位置へ，実際にTD-14で黒石を置いているため，置く場所がなくなったという差のみである．
つまり，CanBlackPlaceはFalseになるはずである．

従って，このテストケースの実行結果が正しいとなると，テストケースID38の状況(境界値分析)において，CalcConditionがBoardのCanBlackPlaceへ与える影響は正しいと言える．

Bに関しては，これが全く反対になり，上記と同様の条件のBoardに対してCalcConditionがBoardのCanWhitePlaceへの作用が正しいと言える
(テストケースIDはそれぞれ，39, 40に対応)．


以上の過程により導出された40個のテストケースを実際にMSTestを用いて実装し，CalcConditionのメソッドテストを行なった．

結果としては，すべて合格となった．

従ってすべてのテストに合格したため，CalcConditionの動作は仕様書から見て正しいと言える．



\section{考察}
\subsection{GameControllerのClickedメソッドをテストするにはどうすれば良いか}
Clickedメソッド入力と，そのメソッドの作用を考え，そこからテストケースを考えたい．

まずClickedメソッドの仕様\cite{journal:sekkeisyo}は以下のとおりである．

\begin{quote}
internal void Clicked(int row, int column) \\

ゲームのメイン処理．MainForm のイベントハンドラから呼び出される．クリックされたセルにピースが置
けるか確認し，置けないならその旨メッセージを表示．置けるなら（Board の Set メソッドで）置いて，
ボードの表示を更新（MainForm の Refresh メソッド）し，AI の処理にうつる．AI はボード上のどこかに置
けるなら置いて（PlaceAiHand メソッド），結果をボードに反映する．AI が置けないなら，AI のターンを飛
ばす旨表示し，プレイヤーが置けるか確認．プレイヤーも置けないなら，ゲームの終了処理（ShowResult メ
ソッド）をする．プレイヤー，AI の手番を飛ばすかどうかは，Board の CanPlace(1 引数)で判断する．
\end{quote}

Internalメソッドであるため，このメソッドに直接作用してテストするのではなく，動作が確認されているMainFormをドライバとして扱い，動作を確認する．

今回，考慮するべき入力は次のとおりである．

\begin{itemize}
  \item Cell \begin{itemize}
    \item 空白か否か
  \end{itemize}
  \item Board \begin{itemize}
    \item ゲームが進行中/終了後の盤面
    \item プレイヤーが置ける/置けない盤面
    \item AIが置ける/置けない盤面
  \end{itemize}
  \item Row, Column
\end{itemize}

また，Clickedの観測可能な動作(出力)は次のとおりである．

\begin{itemize}
  \item MainForm label \begin{itemize}
    \item 置けない旨のメッセージ : A
    \item プレイヤーが置いた旨のメッセージ : B
    \item プレイヤーのターンが飛ばされた旨のメッセージ :C
    \item AIが置いた旨のメッセージ : D
    \item AIのターンが飛ばされた旨のメッセージ : E
  \end{itemize}
\end{itemize}

今回は，A-Eまでの出力に対して，個別にテストを行い，それぞれの独立した動作を検証する．

\paragraph{A}
置けないという事象は，Boardが終了済み，Cellが埋まっている，Row, Columnが異常値である場合に発生する．

従って，デシジョンテーブルは以下の表\ref{tab:enshu6-a}のようになる．

\begin{table}[htbp]
  \centering
  \caption{Aのデシジョンテーブル}
  \label{tab:enshu6-a}
  \begin{tabular}{cccc}
    % ここにカラムを入力
    Cell & Board & Row, Column & A \\
    \hline
    % ここに各セルのデータを入力
    埋まっている & x & x & 表示 \\
    x & 終了済み & x & 表示 \\
    x & x & $Row, Column \leq -1, 8 \leq Row, Column$ & 表示 \\
    \multicolumn{3}{c}{それ以外} & 表示されない \\
  \end{tabular}
\end{table}

xはDontCareである．

以上のデシジョンテーブルをを満たすように，Cell, Board, Row, Columnの状態を同値分割し，テストを行うことで，Aの表示動作が正しいか検証することができる．

同様の手順でBからEに対しても，その出力が動作するタイミングに必要なCell, Board, Row, Columnを同値分割し，それぞれに対してテストを行うことができれば
Clickedメソッドの動作をテストできる．

しかし，今回の問題としては，A-Eはいずれも排他的ではなく，同時に状況が発生する場合がある．
その際に，どの表示が出るのが正しいのかを検証するために，A-Eいずれか複数の発生条件が同時に満たされる場合を考えテストケースを作成する必要があることにも留意したい．

今回は具体的なテストケースの掲載は避ける．

\newpage
\subsection{演習6 テストファーストと継続的インテグレーションについて調べ，関係性をまとめよ\cite{wiki:aws-cicd-definition}, \cite{wiki:tdd}}
まず，テストファーストとは，プロダクト開発において，要件を満たすテストを設計し，設計されたテストを満たすようにプログラムを記述するという考え方である．
特に，テストファーストは近年テスト駆動開発(TDD)として，フレームワーク化され，活用が進んでいる．

特にプログラミングの最終目的が，テストを満たすことになるおかげで，テストの自動化や，成果物の品質を保証しやすいといった点が特徴である．

次に，継続的インテグレーション(CI)とは，プログラミングされたプログラムをデプロイするまでの一定の流れを自動化し，品質向上や工数削減，提供速度の向上などを図る手法である．
具体的に，プログラムのデプロイまでには，プログラムのマージ，ビルド(コンパイル)，テストなどの作業が必要となる．
この一連の流れを，自動化サイクルとして実行し，効率化を図るのがCIということになる．
また，ビルド方法の属人化やテスト工程の複雑化によるヒューマンエラーの削減など，様々な点において，メリットが多い．
さらにこれらの作業をコードベースで自動化できるということは，IaCなどの面から見ても非常に効率が高い．

CIはCI/CDの文脈で利用されることが多く，テスト完了したプログラムをそのままDelivery，つまり完全にデプロイし，クライアントに配信するまでを一連の流れで行うことがおおい．

特に，TDDと相性がいい．
なぜなら，そもそもTDDはテストを満たすプログラムを記述することであり，CIのゴールと一致しているためである．

特に実際の開発においては，個人的なTDDとコードのバージョン管理ツール上での外部TDDを行うことが重要となり，外部TDDにおいてCIを採用するというのがベストプラクティスとなる．

\bibliographystyle{jplain}
\bibliography{Uwano-05-34-MasuiShu.bib}
\end{document}
